\chapter{RISK ANALYSIS AND MITIGATION}%
\label{chap:risk_analysis_and_mitigation}

Risk analysis is the process of identifying potential issues that may affect the successful 
completion of a project and defining strategies to mitigate their impact. In embedded system 
projects, risks may arise from both hardware and software components.

Table~\ref{tab:risk_analysis} summarizes the identified risks associated with the \textbf{\Project{}} 
project along with their mitigation strategies.

\section{RISK, IMPACT AND MITIGATION STRATEGY}

\begin{table}
\centering
\begin{tabular}{|p{4cm}|p{4cm}|p{6.5cm}|}
\hline
\textbf{Risk} & \textbf{Impact} & \textbf{Mitigation Strategy} \\
\hline
Incorrect timer configuration & Inaccurate time display & Careful calculation of timer values and verification using simulation \\
\hline
Display flickering due to improper multiplexing & Poor visual output & Adjusting refresh delay and ensuring correct digit enable logic \\
\hline
Internal oscillator inaccuracy & Time drift over long duration & Acceptable tolerance for academic project and proper prescaler selection \\
\hline
Hardware wiring errors & System malfunction & Verification of circuit connections using Proteus before testing \\
\hline
Software bugs or logic errors & Unexpected system behavior & Incremental testing and use of interrupt-driven design \\
\hline
\end{tabular}
\caption{Risk Analysis and Mitigation Strategies}%
\label{tab:risk_analysis}
\end{table}

\section{SUMMARY}
The identified risks were effectively managed through careful design, simulation-based testing, 
and incremental debugging. As a result, the project was completed successfully with reliable 
performance.
